\chapter{COLONNELLO}\label{colonnello}

\hfill\break

L'Unef fece arrivare una mezza dozzina di Poiane un paio d'ore dopo, nel
frattempo noi avevamo cercato sopravvissuti e assistito i feriti. Con
mio sollievo, trovammo un capitano dell'esercito che prese il comando e
io potei tornare a eseguire ordini.

Tre giorni dopo, al mattino, ci fu un violento temporale mentre stavo
lavorando in una squadra di pulizia; dopo che la pioggia si fu fermata e
il sole fu uscito a sollevare nuvole di vapore dalla pista del campo
d'aviazione, un soldato mi venne a cercare e mi disse che avevo l'ordine
di fare rapporto al maggiore Perkins all'ufficio amministrativo. Il
maggiore Perkins? Che diavolo ci faceva qui, mi aveva detto che era
stata assegnata al settore indiano. Pensai che non potesse significare
nulla di buono per me.

«Sergente Bishop a rapporto, signora.» Ero un po' a corto di fiato per
avere salito le scale di corsa.

Il maggiore Perkins mi guardò per un istante con sorpresa. «Non potevi
ripulirti? Che cosa stavi facendo?»

Abbassai lo sguardo sui miei vestiti sudici e fuligginosi e sulle mie
mani sporche, con le unghie annerite. «Mi è stato detto di venire qui di
corsa, signora. Stavo aiutando a spostare, ehm, detriti dalla pista.»
Detriti che includevano Poiane distrutte, con resti umani ancora a
bordo.

«Oh, diavolo, Bishop. Non c'era mica bisogno che arrivassi fin qui di
corsa, cazzo. Siediti.» Sembrava stanca quanto me. «Come ti senti?»

«Stordito, signora», risposi con sincerità. «Come mi aveva detto, stavo
tenendo un profilo basso qui, al servizio di scorta. Mi facevo gli
affari miei. Poi -- bum! -- si è scatenato l'inferno.» Rabbrividii senza
volerlo, pensando al cratere al posto della sala mensa. «Perché diavolo
hanno colpito la sala mensa, signora? Un paio di minuti di scarto quella
mattina e sarei stato lì.»

Perkins guardò fuori dalla finestra, da cui si vedeva bene quel cratere.
«Il comandante ruhar dell'attacco qui, l'abbiamo interrogata, dice che i
colpi alla mensa, alla caserma principale e all'edificio amministrativo
erano intenzionali, e pensavano che la sala mensa e la caserma sarebbero
stati per lo più vuoti a metà pomeriggio. Se non ci fosse stato il
generale Gupta, la sala mensa sarebbe lo sarebbe stata davvero e non è
possibile che la flotta dei ruhar sapesse del suo arrivo prima di
saltare in orbita. Ha detto che è stata una sfortuna che la sala mensa
fosse piena di gente, che non volevano causare vittime, a meno che non
fosse necessario. Non so se crederle o meno, giudicherà la G2, o i
kristang: la consegneremo a loro domani. Ecco perché i ruhar erano così
incazzati per il fatto che avessi abbattuto quelle due Balene, che
avevano a bordo quasi cinquecento uomini, più gli equipaggi. A quanto
pare, i criceti non si erano accorti di averne uccisi a loro volta la
metà in sala mensa: è un cratere e non avevano perso tempo a esaminarlo
alla ricerca di resti. Per i ruhar, finché non abbiamo detto loro delle
vittime, le vostre azioni avevano inasprito senza motivo il conflitto
qui, in un momento in cui stavano offrendo un cessate il fuoco.»

«Non sembrava così quando le loro cannoniere stavano mitragliando ogni
umano in vista, signora», dissi con fervore.

«Capito, Bishop, tieni a mente che i nostri stavano sparando contro di
loro, quindi diamo la colpa alla nebbia della guerra. Non ho niente
contro le tue azioni, infatti ti ho proposto per un encomio. I nostri
amici di sopra hanno avuto un'idea diversa. Vogliono che l'Unef ti
promuova.» Appariva amareggiata, come se non approvasse quell'idea.

«Primo sergente? Sono stato sergente soltanto per\ldots»

«Bishop, non ti nomineremo primo sergente.» L'espressione sul suo volto
era impossibile da decifrare; pensavo intendesse che i kristang
volessero nominarmi primo sergente, ma l'Unef non l'avrebbe fatto,
perché non ero pronto. Un sentimento con cui ero d'accordo al 100\%, non
ero ancora sicuro di quella cosa del sergente. Forse la divisione aveva
mandato lì Perkins per indorare la pillola, perché avevamo già lavorato
insieme. «I kristang promuovono basandosi quasi soltanto sul successo in
battaglia, ci sono considerazioni politiche e rivalità all'interno dei
clan, ma nel loro sistema le promozioni sono assegnate in virtù del
successo in combattimento, valutato in base al numero di nemici uccisi.
Tra quelle due Balene e le vittime a terra quando le munizioni a bordo
sono esplose, le tue azioni hanno ucciso ben più di mille ruhar. I
kristang sono impressionati. Sono incazzati perché la maggior parte
delle unità umane non ha fatto granché mentre i ruhar erano qui, non
conta se tutti gli effetti non c'era molto che potessimo fare da terra,
con i ruhar che potevano colpire l'intero pianeta dall'orbita. Tu hai
ucciso un bel po' di loro, mentre la maggior parte degli uomini
dell'Unef, specie i superiori, non hanno fatto un bel niente, secondo i
kristang.» Tra i superiori c'era il maggiore Perkins stesso, non era
difficile capire cosa ne pensasse. «I kristang vogliono che noi\ldots{}
ah, cazzo, ecco.»

Prese una piccola scatola da una tasca, la squadrò con disgusto e l'aprì
con un gesto stizzito. All'interno c'erano un paio di mostrine
d'argento, un'aquila che stringeva delle frecce e un ramo d'ulivo negli
artigli, con la testa rivolta verso le frecce. Erano le mostrine
dell'Aquila da guerra. L'esercito americano non le aveva più emesse
dalla seconda guerra mondiale.

Quelle aquile erano le mostrine di un colonnello a tutti gli effetti
nell'esercito, nell'Aeronautica e nei Marine, o di un capitano nella
Marina. «Wow, signora, la stanno promuovendo di due gradi?» Ero colpito,
Perkins era un maggiore e il grado superiore a quello di maggiore era
tenente colonnello. Per quanto ne sapevo, nessuno era mai passato
direttamente dal grado di maggiore a quello di colonnello.

«No, Bishop, idiota», disse Perkins con irritazione. «Te. Sono per te. I
kristang vogliono che l'Unef ti nomini colonnello.»

«Porca puttana.»

«Sì, porca puttana è il minimo. Colonnello è un rango inferiore a quello
che volevano i kristang, si aspettavano che l'Unef ti nominasse
generale. Un dannato generale» Scosse la testa incredula. «Riesci a
credere a una cagata del genere?»

«Signora, non riesco a credere che mi facciano colonnello, figuriamoci
generale.»

«Bishop, non montarti la testa. Sei un soldato abbastanza intelligente,
sei flessibile, ti adatti alle situazioni e hai dimostrato capacità di
prendere iniziative, a volte. Il fatto è che non sei più intelligente,
flessibile o innovativo di quanto l'esercito si aspetti da qualunque
soldato.»

«Sì, signora», concordai, perché era vero.

«Comunque, non sarai generale, abbiamo detto ai kristang che il ruolo
dei generali nei nostri eserciti è per lo più amministrativo, che quello
di colonnello è il più alto grado tra quelli coinvolti in modo diretto
in combattimento. È vero. L'hanno capito, quindi basta che tu sia
colonnello. Quel grado», diede un colpetto alla scatola con le aquile
d'argento, «è abbastanza alto da dimostrare che l'Unef apprezza il
successo in combattimento tanto quanto i kristang. È importante
accontentare i nostri alleati.»

«Ehm», borbottai, ipnotizzato da quelle belle aquile d'argento. Cazzo.
Prima che i ruhar attaccassero, tutto quello che volevo dal mio servizio
militare era pagare il college da qualche parte e andarmene. Ora c'erano
un paio di aquile d'argento sedute di fronte a me. «Non sono bravo a
scrivere lettere e cose del genere.»

«Lettere?»

Staccai i miei occhi dalle aquile per guardarla dritta nei suoi. «Sa,
tipo scrivere una lettera che dice qualcosa come: ``Grazie, è un grande
onore ma non posso accettare, nell'esercito americano non funziona
così\ldots''.»

«Sergente, non riesco a farmi capire, quindi te lo spiego in stile
Barney», sputò con esasperazione e nemmeno un accenno di sorriso, quindi
seppi che non stava usando ``Barney'' in senso ironico. «Questa è
comunque l'ultima volta che posso darti un ordine, prima che tu mi
superi di grado. L'Unef non vuole che rifiuti in modo educato la
promozione. Devi accettare con entusiasmo, accettare la promozione come
tuo diritto per avere ucciso un sacco di ruhar. Devi dire ai kristang
che ti dispiace solo di non aver potuto uccidere più criceti. Sii
fiducioso, audace, assetato di sangue. Sii quello che i kristang
vogliono che gli umani siano, perché il servizio di guarnigione su
Paradiso può essere un lavoro di merda, ma è il lavoro che abbiamo
concordato e i nostri alleati sono il nostro unico passaggio per casa. E
la nostra unica fonte di cibo.»

«Porca puttana», ripetei. Non sapevo cosa dire.

Io. Un colonnello.

Un colonnello a tutti gli effetti.

Io.

«Cosa farò? Come colonnello?» Un colonnello dell'esercito americano
potrebbe essere comandante o vicecomandante di una brigata, ovvero
migliaia di truppe. Non esisteva al mondo che fossi qualificato per
farlo.

Il maggiore Perkins fece un'alzata di spalle ed evitò il mio sguardo:
«Che mi venga un colpo se lo so. Ci penserà l'Unef».

Merda. Mi colpì. L'Unef voleva usarmi come trovata pubblicitaria,
esibirmi ai kristang come esempio del guerriero umano ideale, mentre
alle mie spalle tutti gli umani ridevano di me. Non potei impedire alla
mia faccia di mostrare il mio disgusto.

«Bishop, l'Unef ne ha bisogno. Non deve piacerti, devi fare il tuo
dovere», mi ammonì il maggiore Perkins.

«Certo. Indosserò medaglie sul petto, parlerò con le truppe e forse
venderò dei titoli di guerra. Merda.» Una promozione doveva essere una
buona cosa. «Sarò l'unico burattino che l'Unef sta promuovendo?»

«No, c'è un maggiore dell'esercito cinese che stanno promuovendo tenente
colonnello. I kristang sono rimasti soddisfatti anche delle azioni di un
capitano dell'esercito statunitense, due indiani e uno francese. Sono
tutti morti in battaglia, quindi solo voi due siete vivi per ricevere
l'onore.»

«Questo tizio cinese, viene promosso di un grado soltanto? Perché?»
Cazzo, stavo passando direttamente dal grado di sergente più basso che
ci sia a quello di colonnello.

«La sua unità ha difeso un complesso di magazzini che apparteneva ai
ruhar e l'Unef ora utilizza come deposito di rifornimenti, a quanto pare
i criceti ci avevano lasciato qualcosa di importante quando i kristang
hanno preso il sopravvento. I ruhar volevano a tutti i costi tornare là
dentro. Non potevano rischiare di danneggiare qualsiasi cosa ci fosse
nei magazzini, così hanno fatto sbarcare le truppe e hanno combattuto a
terra. Questo maggiore Chang aveva ancora il controllo di due magazzini
quando la flotta kristang è tornata per inseguire i ruhar, ma ha perso
l'80\% dei suoi uomini. Compreso, si dice in giro, il figlio di un alto
funzionario del governo cinese. Ecco perché sta ottenendo un aumento di
un grado soltanto. Senti, Bishop, cosa sarà dipende tutto da te, chiaro?
Sei un sergente dannatamente bravo, hai fatto un buon lavoro con la tua
squadra integrata di osservazione a Teskor e nessuno dirà che non ti sei
meritato una qualche promozione con le tue azioni qui. Questo è un bene
per l'umanità: dobbiamo indurre i kristang a pensare che le truppe umane
siano preziose.»

\hfill\break

La cerimonia di promozione si svolse a Forte Olimpo, il complesso del
quartier generale dell'Unef. Mi stupì vedere che il posto era
immacolato, i ruhar non l'avevano neanche sfiorato durante l'assalto. Il
maggiore Perkins, che mi accompagnò sul volo in Dumbo per Olimpo, mi
disse che i prigionieri ruhar avevano spiegato all'Unef di non aver
colpito lì perché volevano che il comando dell'Unef restasse intatto, in
modo che qualcuno dotato di autorità potesse ordinare a tutti gli umani
su Paradiso di deporre le armi. Sarebbe stato più facile se i criceti
non avessero disturbato tutte le nostre comunicazioni. Chissà come
pensano gli alieni, eh?

Incontrare il generale Meers e gli altri comandanti anziani non
m'intimidì come mi aspettavo, avevo temuto una lunga cerimonia formale
piena di discorsi, durante la quale sarei stato costretto ad assumere un
minaccioso cipiglio guerriero o un gradevole sorriso, qualsiasi cosa il
funzionario delle pubbliche relazioni dell'Unef avesse ritenuto
opportuna. Invece, poiché in mancanza di tute di protezione complete
solo un piccolo numero di kristang era stato autorizzato a sbarcare su
Paradiso, nessuno di loro era presente a Olimpo. La cerimonia di
promozione, perciò, si svolse nell'ufficio del capitano Meers, che
appuntò le aquile d'argento sulla mia uniforme, mentre una mezza dozzina
di altri superiori stavano a guardare. Dopo la breve cerimonia, ci fu un
delizioso pasto nella mensa degli ufficiali, con bistecche che avevano
un sapore fresco perché, mi dissero, erano state irradiate e raffreddate
invece che congelate, fagiolini che non erano la solita sbobba molle
dell'esercito, patate al forno con burro vero, torta al cioccolato e
autentico caffè appena fatto.

I superiori mangiavano bene. A questo mi potevo abituare. La mia gioia
durò fino a quando un colonnello del corpo dei Marine degli Stati Uniti
mi ringraziò per avere dato all'unità del quartier generale una scusa
per fare una festa; il generale Meers aveva disposto che i superiori a
Forte Olimpo mangiassero razioni da campo sei giorni a settimana, per
ricordare ai comandanti e al loro stato maggiore quello che le truppe
stavano vivendo sul campo. Non sapevo molto di Meers, ma dopo avere
sentito questo guadagnò molti punti ai miei occhi.

La bistecca era ottima. Avevo ancora voglia di un cheeseburger.

\hfill\break

Dopo cena, il generale Meers volle parlare con me in privato. Con me. Il
comandante di tutte le forze umane di Paradiso voleva fare una
chiacchierata con me.

«Serg\ldots{} cazzo, colonnello, è una novità anche per me. Bishop.» Il
mio cognome era opportunamente di rango neutro. «Hai avuto una carriera
pazzesca, davvero interessante, per uno così giovane. Nigeria, poi
l'incidente ruhar nella tua città natale e la cattura di uno di loro. La
mia G2 mi ha detto che sei stato la nostra fonte di informazioni per un
po'. E poi metti insieme una squadra per conto tuo e abbatti due Balene.
Sono rimasto di sasso quando mi hanno detto che sei lo stesso Bishop che
ha catturato un soldato ruhar con un furgone dei gelati. Sono delle gran
cazzo di coincidenze.»

«Mia madre diceva che ero una calamita per i problemi, da ragazzino
riuscivo sempre a cacciarmi in un guaio o nell'altro.» Potevo sentire la
sua voce al mio orecchio mentre lo dicevo. «Non è poi una gran
coincidenza, signore. Ero in congedo quando i ruhar si sono schiantati
nella mia città natale, è stata fortuna. Ma non è stata una coincidenza
che la Bürgermeister\ldots»

«La chi?»

«Il governatore regionale ruhar, o qualunque cosa sia. La chiamavamo
Bürgermeister, quella che mi dava informazioni sui wormhole e tutto il
resto. Non è una coincidenza che abbia scelto me per parlare: mi ha
cercato perché il criceto che ho catturato a Thompson Corners aveva
riferito di essere stato trattato bene, quindi lei voleva incontrarmi.
Mi trovavo a Forte Freccia perché, ehm, mi hanno detto che era perché i
kristang stavano ficcando il naso in giro per cercare di scoprire chi ci
fornisse informazioni, perciò non è stata una coincidenza che io fossi
lì quando i ruhar hanno attaccato. È probabile che sarei stato nella
sala mensa di Forte Freccia con tutti gli altri, non fosse che il
capitano Price ha detto di essere stanco del fatto che l'Unef usi Forte
Freccia come discarica e mi ha messo in servizio di scorta per togliermi
di mezzo. È per questo che non ero alla base quando i ruhar ci hanno
colpito.» Era andata più o meno così, il generale Meers non aveva
bisogno dei dettagli.

«Una cosa tira l'altra, uhm?»

«È così che la vedo, signore.»

«Suppongo che tu abbia ragione su questo. Ora che sei qui, non
aspettarti di poter fare il furbo perché hai l'Unef in pugno. I kristang
assegnano promozioni in virtù del successo in combattimento, quello che
forse non sai è che possono dare il benservito alla gente altrettanto in
fretta.»

«Signore, non so cosa farò adesso, ma qualsiasi compito mi venga
assegnato, lo svolgerò meglio che posso.»

«Bene. L'ultima cosa di cui abbiamo bisogno è che tutti nella cazzo di
Unef pensino di poter svoltare in fretta se fanno qualcosa di
spettacolare. La maggior parte dei tenenti là fuori pensano di essere
più intelligenti dei loro comandanti. E i sergenti pensano di essere
soldati migliori di qualsiasi tenente in pantaloni eleganti.» Sbuffò.
«Diavolo, è probabile che i sergenti abbiano ragione. Abbiamo già avuto
incidenti di persone che cercano di fare cagate alla Rambo fuori di
testa per farsi notare. Quindi, retrocederti al rango di sergente o di
soldato perché hai fatto qualche casino mi renderebbe la vita più
facile, chiaro?»

«Sì, signore.» Era un po' quello che mi aspettavo che sarebbe successo
alla fine. Non avevo il diritto di indossare aquile d'argento.

«Detto questo», Meers soppesò la visuale fuori dalla finestra, dove una
coppia di Polli stava passando in volo, «non voglio che tu fallisca.
Bishop, ti sto mettendo in guardia sul fatto che abbiamo intenzione di
metterti sotto come un mulo a noleggio, quindi se pensi che essere un
colonnello mustang significhi volare in giro e fare discorsi, è meglio
che ti levi subito quel pensiero dalla testa.»

«Mai pensato, signore.» Apprezzai l'avvertimento.

\hfill\break

Il giorno dopo iniziai un corso intensivo sulle responsabilità di un
colonnello dell'esercito americano e sui protocolli e l'etichetta per
trattare con i kristang. Uno dei protocolli per un incontro con loro
prevedeva di sottoporsi a una dieta insipida a partire dal giorno prima;
i kristang ritenevano che gli umani avessero un cattivo odore e che noi,
che ci ingozzavamo di carne o cibi piccanti, avessimo un odore ancora
peggiore. L'Unef aveva sentito dai ruhar che i kristang pensavano che
tutte le altre specie puzzassero, quindi gli umani non avrebbero dovuto
prenderla sul personale. Questo significava fare colazione con farina
d'avena e tè, anziché uova e caffè; pranzo e cena furono altrettanto
insipidi e noiosi. Stando ad altre istruzioni, non si doveva provare a
stringere la mano ai kristang, non gradivano affatto che li toccassimo.
E non si doveva parlare, a meno che i kristang non facessero una
domanda. Inoltre, non si doveva sorridere: i kristang interpretavano i
sorrisi degli umani come il segnale che non stavano prendendo le cose
sul serio e le lucertole in generale non erano conosciute come una
specie allegra.

Le istruzioni riguardanti il mio nuovo ruolo di colonnello erano
semplici e complesse insieme. Semplici, perché era chiaro che il
colonnello incaricato di mettermi al corrente della situazione non
sapeva cosa l'Unef avesse in mente di farsene di me, o se il mio grado
fosse una trovata pubblicitaria a breve termine, perciò le sue
istruzioni consistevano nel dirmi di comportarmi come un ufficiale e di
non fare nulla che potesse imbarazzare l'Unef. Complesse, perché il
colonnello mi spedì per e-mail un'enorme serie di documenti da leggere,
a partire dal materiale formativo che un sottotenente avrebbe dovuto
conoscere a memoria. La guerra sarebbe finita prima che avessi avuto la
possibilità di leggere metà della merda nella mia casella postale.

La mattina successiva, dopo un'altra colazione insipida, i kristang
mandarono una navicella a prendere il generale Meers, diversi membri
dello stato maggiore e me. Eravamo fortunati che l'ascensore spaziale, e
quindi la stazione spaziale, fossero alla stessa longitudine di Forte
Olimpo, il che significava che la loro mattina non cadeva nel bel mezzo
della nostra notte. Meers stava per conferire con i kristang in merito a
qualsiasi cosa di cui volessero parlare loro. Stavo partendo anche io,
perciò poteva darsi che i kristang volessero assegnarmi di persona un
premio, o qualcosa del genere. Non avevamo ben chiaro cosa stesse per
accadere, sapevamo soltanto che i kristang avevano chiesto che il
tenente colonnello Chang e io andassimo alla loro stazione spaziale in
cima all'ascensore.

Indossavo un'uniforme da colonnello nuova di zecca, in cui cercavo di
sedermi con attenzione, per non sgualcire i pantaloni dalla piega
impeccabile. Il giorno prima, mi avevano iniettato un'altra dose di
magico rimedio kristang contro la nausea, quindi non avrei vomitato
farina d'avena addosso a me e a tutti gli altri quando la nave sarebbe
stata in orbita a gravità zero. Seduto accanto a me c'era un certo
tenente Reynolds, il cui unico compito era impedirmi di fare o dire
qualcosa di stupido. Tipo farle in automatico il saluto militare.
Insisteva nel farlo lei e nel chiamare me ``signore''. Risultava
innaturale.

Il viaggio in salita andò più liscio di quello da Campo Alfa, almeno per
come lo ricordo: o questa volta avevamo un pilota migliore, o i kristang
si stavano prendendo cura dei loro Vip umani. La prima volta a bordo
della stazione spaziale, tutto quello che avevo visto era l'interno di
un anello di ancoraggio e corridoi usurati quando ci avevano spintonati
dalla nostra nave all'abitacolo dell'ascensore, che era parcheggiato in
fondo alla stazione. La nostra navicella scese in un hangar
d'atterraggio e riuscimmo a vedere la stazione per davvero. Fu
sorprendente, quello che mi ero immaginato dell'architettura d'interni
dei kristang era spoglio e funzionale, qualcosa di industriale e
militare, adatto alla loro casta guerriera. Quello che vidi era per lo
più elegante e funzionale, ma con elementi d'incredibile ricercatezza.
C'erano arazzi appesi alle pareti, che persino io trovai belli, paesaggi
di pianeti lontani, astronavi delineate su una nebulosa, disegni
geometrici intricati, pure dei fiori, oltre alle raffigurazioni dei
guerrieri kristang in battaglia che mi aspettavo. Erano una specie
interessante, così piena di contraddizioni, se quello che mi aveva detto
la Bürgermeister era vero.

La cerimonia non fu granché, marciammo in una grande stanza, con
kristang seduti lungo entrambi i lati e kristang d'alto rango su un
palco rialzato in fondo. Era la prima volta che ne incontravo di persona
ed ero intimidito. Erano tutti più alti della media umana, corpulenti e
muscolosi, e la loro espressione di default sembrava essere un feroce
scherno. Per quanto ne sapevo, era la loro versione di un sorriso
amichevole. Ah, una parola sulle proteste dei kristang contro l'odore
degli umani: per quanto li riguarda, avrebbero potuto usare un
deodorante per ambienti nella stazione. Stando in una stanza, seppur
grande, con un centinaio di kristang o giù di lì, si avvertiva un odore
secco, coriaceo, con un accenno di qualcosa come il sudore di una
giornata.

La mia parte fu un breve discorso scritto per me; il tenente Reynolds mi
aveva detto che i kristang avevano chiesto di rivedere e approvare le
mie osservazioni in anticipo. Il mio intervento, pronunciato in inglese
e tradotto dallo zPhone, era bellicoso e sanguinario come si conveniva.
Quando ebbi fatto un passo indietro, il tenente colonnello Chang fece un
passo in avanti e tenne un breve discorso simile al mio. Mentre Chang
parlava, i miei occhi vagavano tra la folla. La maggior parte dei
kristang sembrava annoiata, anche un po' disgustata. Se ne stavano
seduti con la faccia di pietra, oppure fissavano il soffitto, o
controllavano i messaggi o giocavano con gli zPhone. Riuscivo a provare
simpatia per loro, i kristang seduti lungo i lati della stanza
partecipavano alla cerimonia da volontari obtorto collo, tutto ciò che
volevano era che finisse. Era come trovarsi in una stanza piena di
adolescenti umani. Il fatto di vedere kristang annoiati e indifferenti
mi diceva quanto fossi insignificante, di certo non si curavano di
quello che avevo fatto, né se mi avessero promosso, o se sarei
sopravvissuto al viaggio di ritorno in navicella su Paradiso. Quando
Chang ebbe finito di parlare, un kristang consegnò una scatola al
generale Meers. Lui prese dalla scatola dei nastri d'oro e li attaccò
alle nostre uniformi. Uno dei kristang sul palco si alzò e ci fece il
saluto, che Chang e io ricambiammo. Ed era finita. Meers e gli ufficiali
che avevano davvero peso rimasero a discutere con i kristang, mentre
Chang e io fummo affidati a una lucertola visibilmente risentita di
dover badare a due umili umani.

«Venite con me, inferiori», è quello che disse il traduttore quando il
kristang ci fece cenno di seguirlo. La sua espressione mi confermò
l'accuratezza della traduzione. Senza garbo, ci guidò nel nostro
tortuoso percorso attraverso la stazione, fino a un ponte di
osservazione che era del tutto deserto. C'erano sedie, un paio di tavoli
e grandi finestre con una splendida vista di Paradiso. Fu difficile per
me trattenermi dal gridare: ``Posso vedere la mia casa!''. Perché
pensavo di avere riconosciuto la grande ansa del fiume che si trovava
proprio a sud di Teskor. Il kristang ci invitò con gesto sprezzante a
sederci, mentre lui andava verso quello che sembrava essere un bar
nell'angolo. Era un bar. Tirò fuori un bicchiere, si versò una generosa
porzione di un liquido dorato, aggiunse due cubetti di ghiaccio e si
stravaccò con rabbia su un sedile di fronte a noi: «Dovrei congratularmi
con voi per i vostri successi come guerrieri e darvi il benvenuto nella
nostra gloriosa coalizione. Puah». Tirò fuori la lingua e ci soffiò
contro un lampone.

Lo guardai più da vicino mentre mandava giù un grosso sorso di liquido.
Qualunque cosa fosse, odorava di alcol e mi ricordava un po' la tequila.
Il kristang aveva gli occhi un po' vacui e vitrei.

Merda. Era già ubriaco. Era venuto alla cerimonia ubriaco. Stava bevendo
ancora. Questo non prometteva niente di buono. Chang attirò la mia
attenzione, mentre il kristang guardava fuori dalla finestra. Scossi la
testa. Qualunque cosa fosse accaduta, ci saremmo entrambi comportati nel
modo migliore possibile, in quanto ospiti dei nostri alleati.

«Aaah.» Il kristang sospirò e sprofondò nella sedia, chiudendo gli
occhi. Dentro di me, speravo che si addormentasse e che io potessi
starmene seduto a godermi la visuale in tranquillità, finché qualcuno
non fosse venuto a prenderci: «Il mio capo mi odia. Altrimenti, perché
mi punirebbe al punto da farmi respirare il fetore ripugnante degli
inferiori?».

Immaginando che non si aspettasse una risposta, tenni la bocca chiusa.
Chang fece lo stesso.

Il kristang scosse la testa e ci agitò davanti il suo zPhone: «Mi hanno
detto che i componenti della feccia hanno bisogno di usarli per parlare
tra di loro. Vergognoso. La vostra specie non dovrebbe parlare più
lingue, è un segno di debolezza! Il gruppo dominante sul vostro
disgustoso pianeta avrebbe già dovuto conquistare gli altri. Tu», indicò
me, «vieni dall'Uh-meri-ii-ca?».

«America, sì, signore», risposi. «Sono in servizio nell'esercito
americano.»

«Esercito?» Rise. «Voi inferiori non siete un esercito. Siete bambini
che si trastullano con i giocattoli. Il mio falcongesso domestico
potrebbe uccidere un centinaio di voi. Sono stato costretto a leggere
del vostro patetico pianeta per non far perdere tempo al mio capo. Un
altro segno che mi odia, costringermi a imbrattarmi la mente con la
storia di una specie tanto patetica. La tua A-me-ri-ca aveva armi
nucleari, l'unica nazione ad avere armi nucleari per diversi anni,
eppure non siete stati in grado di usare il vostro vantaggio per
distruggere i vostri nemici e conquistare il vostro pianeta. Siete
deboli e patetici. E vi chiedete perché non abbiamo rispetto della
vostra specie. Tale debolezza mostra una grave mancanza di risoluzione.
Tu», indicò me, «hai servito in combattimento da qualche parte nel tuo
mondo\ldots{} un posto\ldots{} un posto\ldots» Controllò il suo zPhone,
con una mano resa maldestra dall'alcol.

«Nigeria» dissi.

«Nii-gee-rii-ah. Questo posto non ha armi nucleari?»

Scossi la testa, sorpreso: «In Nigeria, no, non hanno testate nucleari».

«Il vostro esercito avrebbe potuto usare armi nucleari senza paura di
ritorsioni, invece i vostri soldati sono stati mandati a inseguire i
nemici attraverso la giungla. Perché? Perché siete deboli e indegni. Se
questi nii-gee-riaa-nih erano un problema così grande, avreste dovuto
sterminarli e prendervi la loro terra.» Si fermò per un altro sorso
della sua bevanda e ci fece un cenno con il bicchiere: «Voi due siete i
migliori guerrieri della vostra specie? Ah! La vostra specie è inutile
allo sforzo bellico; lo sapete? Lasciare che siate voi a pattugliare
questo pianeta significa accontentare un animale domestico viziato e
stupido. Non avremmo mai dovuto portarvi qui! Sapete cosa stanno facendo
gli umani sulla Terra? Invece di lavorare duro per lo sforzo bellico, si
lamentano che stiamo danneggiando l'ambiente sul pianeta. E i vostri
lavoratori si aspettano di ottenere giorni liberi per le vacanze? Gli
schiavi non vanno in vacanza!». Quest'ultima cosa la strillò con furia e
sbatté il bicchiere sul tavolo di fronte a noi. «Ho detto al mio capo
che dovremmo portare altre specie di schiavi sulla Terra e mostrarvi
come servirci nel modo corretto. Voi umani siete pigri e inutili!»
Guardò fuori dalla finestra, Paradiso era sotto di noi. «Voi non ci
dovreste proprio stare laggiù, dovrebbero essere i kristang a occuparsi
dei ruhar! Infidi ruhar, contaminare il nostro pianeta con agenti
biologici, dovremmo ucciderli tutti.» Fumò di rabbia per un minuto,
sorseggiando la bevanda che era quasi finita. «Mi sono offerto
volontario per assumere il siero sperimentale che permetterà ai kristang
di camminare liberamente sul pianeta. Poi vi mostrerò come trattare i
ruhar. Ci sono così tanti ruhar a infestare il nostro mondo, se ne
facciamo fuori qualche migliaio chi ne sentirà la mancanza, eh? I ruhar
sono deboli e rammolliti, ma dar loro la caccia è un bello sport.»
Prosciugò l'ultimo sorso e disse in tono tranquillo: «Non vedo l'ora di
dare la caccia ai ruhar. Sì».

Il kristang, qualunque fosse il suo nome, appoggiò il bicchiere per
terra sul tappeto e si alzò instabile in piedi. «Godetevi la visuale»,
rise e ci salutò con gesto sprezzante, «potrebbe essere la vostra nuova
casa per sempre, come nostri schiavi. O potrebbe diventare la vostra
tomba.» Rise di nuovo e uscì barcollando dalla porta.

«Merda.» Respirai quando sparì dalla visuale.

Chang annuì, poi disse in perfetto inglese: «Penso che questo non sia il
posto giusto per parlare di\ldots{} di quello che il nostro amico là ha
detto».

Senza dubbio, i kristang ci tenevano sotto stretta sorveglianza a bordo
della stazione. «Hai ragione, colonnello Chang.» Mi vergognavo che Chang
parlasse inglese, mentre l'unica parola cinese che io conoscevo era
``pinguah'', che credo volesse dire ``mela''. L'avevo letta su un
biscotto della fortuna e mi era rimasta impressa. Tutta la mia
conoscenza del linguaggio di una delle civiltà più antiche e grandi del
mondo veniva da un biscotto della fortuna. E i biscotti della fortuna
erano americani, non cinesi.

«Tenente colonnello», disse lui in un tono screziato da una vena di
amarezza.

«Ehi, tu eri un vero ufficiale prima di tutto questo. Io indosso
un'uniforme da colonnello e so di essere soltanto una trovata
pubblicitaria per l'Unef.» C'era di sicuro amarezza nella mia voce.
«Sarò felice quando l'Unef troverà una scusa per abbassarmi di nuovo al
rango di sergente e potrò tornare a essere un vero soldato. L'ultima
cosa che voglio è essere un fobbit e stare dietro a una scrivania.»

«Fobbit?» Chang sollevò un sopracciglio.

«Ah, scusa, slang dell'esercito americano. Un fobbit è un tizio che se
ne sta al sicuro dentro la recinzione di una base operativa avanzata,
mentre i veri soldati sono fuori sul campo. Uno che passa le carte
invece che imbracciare un fucile.»

«Ah, sì. Ne abbiamo anche noi nel nostro esercito. Su questo pianeta,
però, non ci sono delle vere e proprie retrovie, penso. Quando il nemico
può attaccare dall'orbita, ovunque sulla superficie del pianeta è prima
linea.»

«Non hai tutti i torti.» Mi alzai e mi diressi verso la finestra. «Ero a
Forte Freccia. Prima, la mia squadra integrata di osservazione si
trovava in un villaggio di criceti accanto all'ansa di quel fiume, a
ovest di quelle montagne», indicai invano. «Dove siete stazionati voi?»

Parlammo delle nostre esperienze su Paradiso e delle nostre vite prima
che i ruhar attaccassero. Chang era stato un ufficiale di artiglieria
nell'Esercito popolare di liberazione, un ufficiale di carriera di una
famiglia militare. Era preoccupato per la sua famiglia sulla Terra
quanto lo ero io: nemmeno i cinesi avevano avuto notizie, niente
messaggi o lettere da casa, niente di niente. Chang era un bravo
ragazzo, un professionista, di sicuro un ufficiale migliore di me. Anche
lui non moriva certo dalla voglia di essere messo in mostra dall'Unef in
funzione di espediente di pubbliche relazioni, anche se pensava che
sarebbe stata una mossa meno importante di quanto mi aspettassi. Durante
l'azione che aveva portato alla sua promozione, aveva perso la maggior
parte dei suoi uomini, tra cui un tenente che era l'unico figlio di un
alto funzionario del governo. L'esercito cinese su Paradiso non ci
teneva a strombazzare troppo l'azione di Chang. Se il figlio di un
senatore fosse stato ucciso tra le mie fila, la mia carriera militare
sarebbe stata di sicuro stroncata.

Il mio zPhone squillò, era il tenente Reynolds che si chiedeva dove
cazzo fossi, solo che lei lo disse in modo gentile. «Siamo su un ponte
di osservazione, dovremmo essere due\ldots», Chang tirò su tre dita,
«no, tre livelli più in basso. C'è un bar qui. Un bar per i kristang»,
mi affrettai ad aggiungere.

«Cos'è successo al kristang che vi è stato assegnato?»

«Aveva, ehm, aveva qualcosa di importante da fare», dissi, consapevole
che i kristang stavano quasi di sicuro ascoltando. Qualcosa di
importante da fare, tipo dormire finché non si sarebbe svegliato coi
postumi della sbornia.

«Signore, per favore, resti dove si trova. C'è il tenente colonnello
Chang con lei?»

«Sì, è qui, non andiamo da nessuna parte.»

Cinque minuti dopo, Reynolds e due kristang incazzati si presentarono
per recuperarci. Ci riportarono all'hangar di atterraggio e, non appena
il generale Meers e il suo stato maggiore furono a bordo, la porta della
navicella si chiuse e suonò l'allarme per depressurizzare l'hangar. A
dispetto della mia paura che fosse in qualche modo incazzato con me,
Meers non disse nulla, Reynolds non sembrava pensare che ci fossero
grossi problemi, e la discesa andò liscia.

Io e Chang non parlammo mai di quello che aveva detto l'ubriacone
kristang, di come avesse definito gli umani ``schiavi''. Ci pensai molto
su e feci rapporto al personale dei servizi segreti di Meers. Con mio
sgomento, si comportarono come se avessero già sentito tutto in
precedenza: non era una notizia e non era una cosa importante. Lo era
per me, confermava le mie peggiori paure: che la Bürgermeister mi avesse
detto la verità sui kristang.

Stavamo combattendo quella guerra dalla parte sbagliata?